\subsection{Критерий устойчивости Стодолы и Рауса--Гурвица}

Корни характеристического уравнения определяют поведение системы управления и, в частности, её устойчивость.

Необходимый (но не достаточный) критерий устойчивости Стодолы заключается в том, что все коэффициенты характеристического полинома должны быть положительными:
\begin{equation}
    a_i > 0,
    \qquad i = 0,1,\dots,6.
\end{equation}

Для исследуемой системы коэффициенты характеристического полинома, полученные в пункте~2.1, имеют вид:
\begin{equation}
\begin{aligned}
    a_6 &= 0.015, \\
    a_5 &= 5.072, \\
    a_4 &= 167, \\
    a_3 &= 3500k_{\text{п}} + 6500, \\
    a_2 &= 14000k_{\text{п}} + 125000, \\
    a_1 &= 3500000k_{\text{п}} + 14000k_{\text{и}}, \\
    a_0 &= 3500000k_{\text{и}}.
\end{aligned}
\end{equation}

Следовательно, необходимое условие устойчивости выполняется при:
\begin{equation}
    k_{\text{п}} > 0,
    \qquad
    k_{\text{и}} > 0.
\end{equation}

Для получения достаточных условий устойчивости воспользуемся критерием Рауса--Гурвица. Сформируем матрицу Гурвица для характеристического полинома шестого порядка:
\begin{equation}
\Gamma =
\begin{pmatrix}
a_5 & a_3 & a_1 & 0   & 0   & 0 \\
a_6 & a_4 & a_2 & a_0 & 0   & 0 \\
0   & a_5 & a_3 & a_1 & 0   & 0 \\
0   & a_6 & a_4 & a_2 & a_0 & 0 \\
0   & 0   & a_5 & a_3 & a_1 & 0 \\
0   & 0   & a_6 & a_4 & a_2 & a_0
\end{pmatrix}.
\end{equation}

Согласно критерию Рауса--Гурвица, система является устойчивой тогда и только тогда, когда все главные угловые миноры матрицы Гурвица положительны:
\begin{equation}
    \Delta_1 > 0,
    \quad
    \Delta_2 > 0,
    \quad
    \dots,
    \quad
    \Delta_6 > 0.
\end{equation}

Для вычисления матрицы Гурвица и её главных угловых миноров использовалась система Wolfram Mathematica.

Исходный код представлен в листинге~\ref{lst:hurwitz}.

\begin{listing}[H]
\caption{Формирование матрицы Гурвица и вычисление миноров}
\vspace{-10pt}
\label{lst:hurwitz}
\begin{minted}[linenos, numbersep=10pt]{mathematica}
a = coeffsDescending;

H = {
 {a[[2]], a[[4]], a[[6]], 0, 0, 0},
 {a[[1]], a[[3]], a[[5]], a[[7]], 0, 0},
 {0, a[[2]], a[[4]], a[[6]], 0, 0},
 {0, a[[1]], a[[3]], a[[5]], a[[7]], 0},
 {0, 0, a[[2]], a[[4]], a[[6]], 0},
 {0, 0, a[[1]], a[[3]], a[[5]], a[[7]]}
};

Delta1 = Det[H[[1 ;; 1, 1 ;; 1]]];
Delta2 = Det[H[[1 ;; 2, 1 ;; 2]]];
Delta3 = Det[H[[1 ;; 3, 1 ;; 3]]];
Delta4 = Det[H[[1 ;; 4, 1 ;; 4]]];
Delta5 = Det[H[[1 ;; 5, 1 ;; 5]]];
Delta6 = Det[H[[1 ;; 6, 1 ;; 6]]];
\end{minted}
\end{listing}

Главные угловые миноры матрицы Гурвица представляют собой
громоздкие полиномиальные выражения относительно параметров
$k_{\text{п}}$ и $k_{\text{и}}$, поэтому их явный вид не приводится.

На основании условий
\[
\Delta_i > 0,
\qquad i = 1,\dots,6,
\]
была построена область устойчивости системы управления в
пространстве параметров $(k_{\text{п}}, k_{\text{и}})$.

Для построения области устойчивости в пространстве параметров $(k_{\text{п}}, k_{\text{и}})$ использовалась функция Хевисайда:
\begin{equation}
    D(k_{\text{п}}, k_{\text{и}})
    =
    \theta(\Delta_1)
    \theta(\Delta_2)
    \theta(\Delta_3)
    \theta(\Delta_4)
    \theta(\Delta_5)
    \theta(\Delta_6),
\end{equation}
где $\theta(x)$ --- функция Хевисайда.

Функция $D(k_{\text{п}},k_{\text{и}})$ принимает значение, равное единице, если все главные угловые миноры матрицы Гурвица положительны, и ноль в противном случае.

В Wolfram Mathematica область устойчивости строилась следующим образом (листинг~\ref{lst:stability_region}):

\begin{listing}[H]
\caption{Построение области устойчивости}
\vspace{-10pt}
\label{lst:stability_region}
\begin{minted}[linenos, numbersep=10pt]{mathematica}
stab[kp_, ki_] :=
 UnitStep[Delta1] *
 UnitStep[Delta2] *
 UnitStep[Delta3] *
 UnitStep[Delta4] *
 UnitStep[Delta5] *
 UnitStep[Delta6];

 RegionPlot[stab[kp, ki] == 1, {kp, -1, 15}, {ki, -1, 45},
  AxesLabel -> {"kп", "kи"}, FrameLabel -> {"kп", "kи"},
  PlotPoints -> 100, MaxRecursion -> 4, ImageSize -> Large,
  LabelStyle -> Directive[Black, 22], BaseStyle -> {FontSize -> 18}]
\end{minted}
\end{listing}

Область устойчивости системы представлена на рисунке~\ref{fig:pdf/stability_region}.

\screenshot[0.55]{pdf/stability_region}{Область устойчивости системы}

Для проверки выберем точку из найденной области устойчивости:
\begin{equation}
    k_{\text{п}} = 10,
    \qquad
    k_{\text{и}} = 15.
\end{equation}

в Wolfram Mathematica были вычислены корни характеристического полинома (рис.~\ref{fig:pdf/roots_wolfram}).

\screenshot[0.95]{roots_wolfram}{Вычисление корней характеристического полинома в Wolfram Mathematica}

Для выбранных значений параметров корни характеристического полинома имеют вид:
\begin{equation}
\begin{aligned}
    p_1 &= -329.53, \\
    p_{2,3} &= -2.54315 \pm 85.4706i, \\
    p_4 &= -1.50409, \\
    p_{5,6} &= -1.00628 \pm 31.0608i.
\end{aligned}
\end{equation}

Так как действительные части всех корней отрицательны, выбранная точка принадлежит области устойчивости.
